% !TeX program = pdflatex
% Self-contained source: pdflatex twice; no fontspec or external bibliography.
\documentclass[11pt,a4paper]{article}
\usepackage[margin=24mm,headheight=14pt]{geometry}
\usepackage[T1]{fontenc}
\usepackage[utf8]{inputenc}
\usepackage{lmodern,microtype}
\usepackage{amsmath,amssymb,mathtools,amsthm}
\usepackage{enumitem,booktabs,longtable,array,xurl,listings,needspace}
\usepackage[hidelinks,unicode]{hyperref}
\hypersetup{pdftitle={Economic Theory: Proof Attempts and Adversarial Checks},pdfauthor={Research notes},pdfsubject={Eleven user-supplied mathematical problems}}
\newcommand{\NE}{\operatorname{NE}}
\DeclareMathOperator*{\argmax}{arg\,max}
\DeclareMathOperator*{\argmin}{arg\,min}
\newtheorem{theorem}{Theorem}[section]
\newtheorem{lemma}[theorem]{Lemma}
\newtheorem{proposition}[theorem]{Proposition}
\newtheorem{corollary}[theorem]{Corollary}
\theoremstyle{definition}
\newtheorem{definition}[theorem]{Definition}
\theoremstyle{remark}
\newtheorem*{remark}{Remark}
\setlength{\parindent}{0pt}
\setlength{\parskip}{5pt plus 1pt minus 1pt}
\setlength{\emergencystretch}{2em}
\setcounter{tocdepth}{1}
\setlist{leftmargin=*,itemsep=2pt,topsep=4pt}
\lstset{basicstyle=\small\ttfamily,columns=fullflexible,keepspaces=true,
 breaklines=true,breakatwhitespace=false,showstringspaces=false,
 frame=single,framerule=0.25pt,aboveskip=8pt,belowskip=8pt,xleftmargin=3pt,xrightmargin=3pt}
\allowdisplaybreaks[1]
\raggedbottom

\begin{document}
{\large\bfseries Research attempt and adversarial verification}\par
9 October 2026\hfill Original-target status: \textbf{UNRESOLVED}\par
\section{C73-15: Payoff characterization for a fixed mirror-descent rule}
\label{sec:C73-15}
\noindent\textbf{Input:} \texttt{C73-15.tex}. \quad\textbf{Tools:} web browsing [YES]; Python computation [YES].

\subsection*{1) FORMAL RESTATEMENT}
Fix finite nonempty action sets, proper lower-semicontinuous strictly convex
regularizers $h_i$ on their simplexes, differentiable and finite on the relative
interiors, and a sequence $\eta_t>0$. Let $X=\prod_i\Delta(S_i)$ and
$v_i(x)_a=u_i(a,x_{-i})$. For an interior starting point define
\[
 D_{h_i}(p,x)=h_i(p)-h_i(x)-\langle\nabla h_i(x),p-x\rangle,
\quad
 x_i^{t+1}=\argmax_{p\in\Delta(S_i)}
 \{\eta_t\langle v_i(x^t),p\rangle-D_{h_i}(p,x_i^t)\}.
\]
As in the input, assume uniqueness of every update and that all iterates remain
where the required gradients are defined. The target is a condition on $u$ for
\[
 \forall x^1\in\operatorname{relint}X,\qquad
 \operatorname{dist}(x^t,\NE(G))\longrightarrow0.
\]
The rule, including every step size, is fixed before testing a game. The input
does not select one particular regularizer or schedule, so a theorem for a
special choice is a partial classification, not a solution for the whole family.

The main specialization below is negative entropy
$h_i(p)=\sum_a p_a\log p_a$, with $0\log0=0$ and natural logarithms. It gives
\[
 x_{i,a}^{t+1}=
 \frac{x_{i,a}^{t}\exp(\eta_t v_{i,a}(x^t))}
      {\sum_bx_{i,b}^{t}\exp(\eta_t v_{i,b}(x^t))}.
 \tag{MD}
\]
The formula is the unique optimizer: the first-order conditions of the strictly
concave entropy-regularized objective give the displayed positive probabilities,
and normalization determines the multiplier. Hence the interior-domain condition
is satisfied for every finite step.

\subsection*{2) QUICK LITERATURE/CONTEXT CHECK}
Bichler et al. \cite{potentialness} distinguish potential structure from general
learning convergence and explicitly discuss the absence of a comprehensive
classification. Their empirical potentialness measure is not a necessary-and-sufficient
theorem for the particular discrete rule above. The following schedule-sensitive
statements and counterexamples are proved directly.

\subsection*{3) ATTACK PLAN}
\textbf{Proof track.} Obtain exact log-ratio recursions for entropy; classify the
finite-total-step case using initializations near every vertex; identify positive
subclasses when total learning time is infinite.

\textbf{Construction track.} Try a one-player game with summable steps, a zero-sum
mixed-equilibrium game, and an identical-payoff game with a large constant step.
These test three different failed generalizations.

\textbf{Tools and purpose.} Entropy first-order conditions make updates explicit;
log odds remove normalization; uniform payoff bounds control finite total movement;
closedness of the Nash set permits vertex tests; strict convexity supplies an
energy inequality; invariant diagonal profiles reduce a two-player update to one
variable; exact logarithmic identities certify a two-cycle.

\subsection*{4) WORK}
\begin{theorem}[Complete classification for summable entropic steps]
For the entropy rule (MD), suppose $S=\sum_{t\geq1}\eta_t<\infty$. Then the universal
convergence property holds if and only if the game is strategically constant:
\[
 u_i(a_i,s_{-i})=u_i(b_i,s_{-i})
 \quad\text{for all }i,a_i,b_i,s_{-i}.
 \tag{SC}
\]
Equivalently, $\NE(G)=X$.
\end{theorem}
\begin{proof}
If (SC) holds, all pure actions of a player have the same payoff against every
mixed opponents' profile. Every mixed profile is therefore Nash, proving sufficiency.

For necessity choose $M$ with $|u_i(s)|\leq M$ for every player and pure profile.
Multilinearity gives $|v_{i,a}(x)|\leq M$ on $X$. Iterating the ratio form of (MD),
\[
 \frac{x_{i,a}^{t}}{x_{i,b}^{t}}
 =\frac{x_{i,a}^{1}}{x_{i,b}^{1}}
   \exp\left(\sum_{s<t}\eta_s[v_{i,a}(x^s)-v_{i,b}(x^s)]\right).
\]
Fix an arbitrary pure profile $a^*$. Take interior starting points with
$x_{i,a_i^*}^{1}\geq1-\delta$. Summing the ratios to $a_i^*$ gives, for every $t$,
\[
 \frac{1-x_{i,a_i^*}^{t}}{x_{i,a_i^*}^{t}}
 \leq e^{2MS}\frac{\delta}{1-\delta}.
\]
Thus every iterate is within a quantity tending to zero with $\delta$ of the
vertex $a^*$. If the assumed universal Nash-set convergence holds, compactness
provides a subsequential Nash limit with the same closeness bound. Letting
$\delta\downarrow0$ and using that the Nash set is closed shows $a^*\in\NE(G)$.
Every pure profile is therefore Nash. For fixed $i,s_{-i}$, the profiles
$(a_i,s_{-i})$ and $(b_i,s_{-i})$ give both opposite best-response inequalities,
forcing equality of the two payoffs. This proves (SC).
\end{proof}
The necessity argument uses the universal interior-initialization quantifier.
It does not assert that one preselected interior starting point must fail in every
nonconstant game.

\begin{proposition}[Two positive tests when total learning time is infinite]
Under (MD), if $\sum_t\eta_t=\infty$, every one-player finite game converges to its
optimal-action simplex. Also, every finite game in which each player has a strictly
dominant pure action converges to the corresponding pure Nash equilibrium.
\end{proposition}
\begin{proof}
For one player, let $c_a$ denote its fixed payoff and
$S_t=\sum_{s<t}\eta_s$. The exact formula is
\[
 x_a^t=\frac{x_a^1e^{c_aS_t}}{\sum_bx_b^1e^{c_bS_t}}.
\]
If $c_a<\max_b c_b$, its ratio to a maximizing action tends to zero. Ratios between
maximizing actions remain their initial ratios, proving the asserted convergence.

If $a_i^*$ strictly dominates every other action, finiteness gives
$v_{i,a_i^*}(x)-v_{i,b}(x)\geq\delta_i>0$ for every $x$ and $b\ne a_i^*$.
The same ratio recursion bounds
$x_{i,b}^t/x_{i,a_i^*}^t$ by its initial value times $e^{-\delta_iS_t}$.
All the undesired masses tend to zero. The dominant profile is Nash by definition.
\end{proof}

\begin{theorem}[Zero-sum does not guarantee last-iterate convergence]
For matching pennies
\[
 A=\begin{pmatrix}1&-1\\-1&1\end{pmatrix},\qquad B=-A,
\]
the simultaneous entropy rule fails to approach the Nash set from every interior
initialization other than the equilibrium, for every positive step sequence.
\end{theorem}
\begin{proof}
Write $p,q$ for the first-action probabilities and
$z=\log(p/(1-p))$, $w=\log(q/(1-q))$. The unique Nash equilibrium is
$p=q=1/2$: a nonindifferent opponent would force a pure best reply, and no pure
pair is mutually optimal; indifference requires exactly these probabilities.
The ratio recursion gives
\[
 z^{t+1}=z^t+2\eta_t\tanh(w^t/2),\qquad
 w^{t+1}=w^t-2\eta_t\tanh(z^t/2).
\]
Set $E(z,w)=\log\cosh(z/2)+\log\cosh(w/2)$. Its Hessian is diagonal with strictly
positive entries $\tfrac14\operatorname{sech}^2(z/2)$ and
$\tfrac14\operatorname{sech}^2(w/2)$. The displacement above is orthogonal to
$\nabla E$. Strict convexity therefore gives
$E(z^{t+1},w^{t+1})>E(z^t,w^t)$ unless $(z^t,w^t)=(0,0)$.
Starting elsewhere, energy stays at least $E(z^1,w^1)>0$. Convergence to the unique
Nash profile would make $(z^t,w^t)\to(0,0)$ and the energy tend to zero, which is
impossible. This argument uses no lower bound on $\eta_t$ and no assumption on
its total sum.
\end{proof}

\begin{theorem}[An exact two-cycle in an exact potential game]
Let both payoff matrices be
$A=B=\left(\begin{smallmatrix}0&1\\1&0\end{smallmatrix}\right)$, take
$\eta_t=4\log3$, and initialize $p^1=q^1=3/4$. Then the entropy iterates alternate
between $(3/4,3/4)$ and $(1/4,1/4)$. Neither profile is Nash.
\end{theorem}
\begin{proof}
The common payoff itself is an exact potential. Symmetry preserves $p=q$.
For the common log odds $z$, the payoff difference between the first and second
actions is $1-2p=-\tanh(z/2)$, so
\[
 z^{t+1}=z^t-4\log3\,\tanh(z^t/2).
\]
At $z=\log3$, $\tanh(z/2)=1/2$, and the next value is $-\log3$.
Oddness gives the reverse step. At $(3/4,3/4)$ the expected payoff is $3/8$ and
switching to the second pure action gives $3/4$; at $(1/4,1/4)$ switching to the
first gives the same improvement. The periodic orbit is disjoint from the closed
Nash set, so its distance from that set is positive.
\end{proof}
This is a counterexample to arbitrary-step potential convergence, not to the
summable-step classification, whose assumptions it does not satisfy.

\subsection*{5) VERIFICATION}
For a one-player two-action game with rewards $(1,0)$, initialization $(1/2,1/2)$,
and $\eta_t=2^{-t}$, the exact limit is
$(e/(1+e),1/(1+e))$, not an optimum. This tiny case checks the finite-total-step
obstruction without interactions. A one-action player contributes no condition
to (SC), and a wholly constant game trivially passes every rule.

The script verifies the potential two-cycle with a floating-point round-trip
error below $10^{-12}$; its proof is the exact identity involving $\log3$ above.
It also checks 100 strict matching-pennies energy increases from log odds
$(0.7,-0.4)$ with step $0.1$. Minimal pseudocode is
\begin{lstlisting}
for t:
    for each player and action:
        new_log_weight = log(old_weight) + eta[t] * payoff_vector[action]
    normalize each player's new weights simultaneously
\end{lstlisting}
Simultaneity matters: using one player's new probabilities in the other player's
update would be a different algorithm. The proof never treats negative entropy
as differentiable on the boundary. None of the failures concerns convergence of
time-averaged play to a correlated equilibrium.

\Needspace{8\baselineskip}
\subsection*{6) FINAL}
\textbf{LABEL: UNRESOLVED.}

\textbf{Strongest checked partial result.} A complete payoff-array classification for
entropy with summable positive steps; positive infinite-total-step subclasses; an
all-positive-schedules matching-pennies obstruction; and an exact potential-game
two-cycle for a specified constant step.

\textbf{Exact first gap.} For arbitrary fixed regularizers and nonsummable schedules,
no necessary-and-sufficient payoff-array criterion is proved.

\textbf{Next three moves.}
\begin{enumerate}[nosep]
\item For nonsummable entropy steps, identify a payoff condition that controls
last-iterate rather than averaged regret and survives all interior initializations.
\item For a specified potential array, determine a sharp step-size stability bound;
the exact cycle supplies a concrete large-step test.
\item Extend the finite-total-step vertex argument to regularizers for which a
uniform, initialization-sensitive movement bound can actually be proved.
\end{enumerate}
\textbf{Counterexample perspective.} Two actions and two players already suffice to
break both zero-sum and potential-based universal claims. They do not establish
that an arbitrary specified rule has no useful convergent class.

\paragraph{Reproducibility and scope.} This is one of eleven companion reports. The bundle includes \texttt{verify.py} and its executed results. Finite experiments supplement the proofs; they are not proofs of asymptotic claims. Literature coverage is targeted, and no novelty or formal proof-assistant certification is claimed.

\begin{thebibliography}{99}
\small
\bibitem{potentialness}
Martin Bichler, Davide Legacci, Panayotis Mertikopoulos, Matthias Oberlechner, and Bary Pradelski.
\emph{Characterizing the Convergence of Game Dynamics via Potentialness}.
arXiv:2503.16285, version 1 (2025).
\url{https://arxiv.org/html/2503.16285v1}.

\end{thebibliography}

\end{document}
